\chapter{The Teaching Mode}
\label{ch:teaching}

\begin{versebox}
Appeal and danger, escape and fruit and means,\\
plus the Buddha's direct command ---\\
that's the Teaching Mode.
\end{versebox}

\noindent Now the deep dive begins. We've seen the overview, the enumeration, the verse summaries. From here on, Mahākaccāyana takes each mode and shows us exactly how it works --- not in theory, but on real suttas, with real examples, line by line.\footnote{This begins the \textit{Paṭiniddesavāro} --- the ``detailed exposition section'' (literally ``counter-description''). Where the Niddesavāro gave compressed verse definitions, the Paṭiniddesavāro unpacks each one in full prose with worked examples. The section heading in the Pali is \textit{Desanāhāravibhaṅgo} --- ``Analysis of the Teaching Mode.'' The word \textit{vibhaṅga} (analysis, breaking apart) is the same word used for the Vibhaṅga texts of the Abhidhamma --- a formal, systematic unpacking.}

And the first mode he picks? Teaching. The most fundamental one. The structure that every lesson the Buddha ever gave follows.

\section*{What the Teaching Mode Teaches}

``What is the Teaching Mode?''\footnote{Pali: \textit{``Tattha katamo desanāhāro? `Assādādīnavatā'ti gāthā ayaṃ desanāhāro.''} --- ``What is the Teaching Mode? The verse `Appeal and danger\ldots' --- that is the Teaching Mode.'' Mahākaccāyana first points back to his own mnemonic verse from Chapter~3, anchoring the detailed analysis to the compressed definition.} Straight to it. And the answer is a list of six things that every teaching contains:

\begin{enumerate}[itemsep=4pt]
\item The \textbf{appeal} of the thing being discussed
\item The \textbf{danger} in it
\item The \textbf{escape} from it
\item The \textbf{fruit} --- the result of escaping
\item The \textbf{means} --- the method for doing it
\item The \textbf{command} --- the Buddha's direct instruction to practitioners
\end{enumerate}

Then the Netti quotes the Buddha's own description of his teaching method: ``I will teach you a Dhamma that is good in the beginning, good in the middle, good in the end --- with meaning and with phrasing, complete and pure --- I will make clear the spiritual life.''\footnote{Pali: \textit{``Dhammaṃ vo, bhikkhave, desessāmi ādikalyāṇaṃ majjhekalyāṇaṃ pariyosānakalyāṇaṃ sātthaṃ sabyañjanaṃ kevalaparipuṇṇaṃ parisuddhaṃ brahmacariyaṃ pakāsessāmī''ti.} This is one of the most frequent stock phrases in the Canon --- it appears at the opening of hundreds of suttas. The Netti takes it as programmatic: ``good in the beginning'' means the teaching opens well (for quick learners), ``good in the middle'' means it develops well (for those who need elaboration), and ``good in the end'' means it concludes well (for those who need the full treatment). ``With meaning'' (\textit{sātthaṃ}) refers to the six meaning-terms; ``with phrasing'' (\textit{sabyañjanaṃ}) refers to the six letter-terms. Even this stock phrase, the Netti shows, is a compressed description of the Buddha's entire pedagogical architecture.}

That phrase --- ``good in the beginning, good in the middle, good in the end'' --- is one you've probably seen a hundred times in the suttas without thinking much about it. The Netti says: actually, this phrase describes a three-stage teaching system designed for three different types of student. We'll get to that in a moment.

But first, let's see how the six components work in practice.

\section*{A Worked Example: The Verses on Sense Pleasure}

The Netti demonstrates the Teaching Mode using a set of verses from the Suttanipāta about sense pleasure.\footnote{These verses are from the \textit{Kāmasutta} (Sn 4.1, verses 766--771), one of the most celebrated poems in the Pali Canon. The Netti uses them because they perfectly demonstrate all six components of the Teaching Mode in sequence. This is typical of the Netti's method: it doesn't invent examples --- it shows you structures that are already there in the Canon, hidden in plain sight.} Watch how the six-part structure emerges:

\subsection*{The Appeal}

\begin{versebox}
When you want something and you get it ---\\
oh, the pleasure! A person is delighted\\
when they obtain what they desired.
\end{versebox}

That's the appeal.\footnote{Pali: \textit{``Kāmaṃ kāmayamānassa, tassa cetaṃ samijjhati; / Addhā pītimano hoti, laddhā macco yadicchatī''ti.} --- ``When a person desiring a pleasure succeeds in getting it, certainly they are delighted in mind, having obtained what they wished for.'' (Sn 4.1, v. 766). The Netti simply quotes the verse and says: \textit{``Ayaṃ assādo''} --- ``This is the appeal.'' No elaboration needed. The verse speaks for itself. \textit{Assāda} (gratification, appeal) is one of the most important terms in the Buddha's analytical vocabulary --- it names the genuine pleasure that comes from getting what you want, which the Buddha never denies exists.} The Buddha doesn't pretend that getting what you want isn't pleasant. It is. He starts there.

\subsection*{The Danger}

\begin{versebox}
But if those pleasures slip away\\
from that person, that desiring creature ---\\
they suffer like one pierced by an arrow.
\end{versebox}

That's the danger.\footnote{Pali: \textit{``Tassa ce kāmayānassa, chandajātassa jantuno; / Te kāmā parihāyanti, sallaviddhova ruppatī''ti.} --- ``But if those pleasures decline for that person, that creature in whom desire has arisen, they suffer like one pierced by a dart.'' (Sn 4.1, v. 767). \textit{Ādīnava} (danger, drawback) is the counterweight to \textit{assāda}. Note the imagery: the arrow (\textit{salla}) suggests something that penetrates and stays embedded. The pleasure of getting what you want is real, but the suffering when you lose it is like a wound that doesn't heal on its own.}

Notice the honesty. He doesn't jump from ``pleasure is bad'' to ``therefore renounce.'' He says: the pleasure is real. \textit{And} --- the suffering when it ends is real too.

\subsection*{The Escape}

\begin{versebox}
Whoever avoids sense pleasures\\
the way you'd avoid a snake's head with your foot ---\\
that person, mindful, goes beyond\\
this sticky trap in the world.
\end{versebox}

That's the escape.\footnote{Pali: \textit{``Yo kāme parivajjeti, sappasseva padā siro; / Somaṃ visattikaṃ loke, sato samativattatī''ti.} --- ``Whoever avoids pleasures as one avoids a snake's head with the foot --- that one, mindful, overcomes this stickiness in the world.'' (Sn 4.1, v. 768). \textit{Nissaraṇa} (escape) completes the triad. The simile is vivid: you don't fight the snake, you \textit{step around it}. And the key word is \textit{sato} --- ``mindful.'' The escape isn't through willpower or suppression; it's through mindfulness. The word \textit{visattikā} (stickiness, attachment) literally means ``that which clings'' --- sense pleasure is portrayed not as evil but as \textit{adhesive}.}

Then the Netti gives a second set of examples --- another round of appeal, danger, and escape --- this time about material possessions.\footnote{The second set uses Sn 4.1, verses 769--771: ``Fields, land, gold, cattle and horses, servants and dependents, women, relatives --- a person craves these various pleasures'' (appeal). ``Weakness overpowers them, troubles crush them --- suffering follows like water flooding a broken ship'' (danger). ``Therefore let a person, always mindful, avoid sense pleasures; having abandoned them, cross the flood, like one who has bailed out the boat and reached the far shore'' (escape). The Netti quotes both sets to show that the pattern \textit{repeats}: the six-part structure isn't occasional --- it's systematic.}

\subsection*{The Fruit}

\begin{versebox}
The Dhamma protects the one who practices it\\
like a great umbrella in the rains.\\
This is the benefit of Dhamma well-practiced:\\
the practitioner doesn't go to a bad destination.
\end{versebox}

That's the fruit.\footnote{Pali: \textit{``Dhammo have rakkhati dhammacāriṃ, chattaṃ mahantaṃ yatha vassakāle; / Esānisaṃso dhamme suciṇṇe, na duggatiṃ gacchati dhammacārī''ti.} The umbrella simile (\textit{chattaṃ mahantaṃ} --- ``a great umbrella'') is both practical and beautiful. In India's monsoon season, a good umbrella is the difference between misery and comfort. The Dhamma does the same: it doesn't stop the storms of life, but it keeps you from drowning in them.} The result of practicing the escape: protection. Not in some abstract philosophical sense --- protection the way an umbrella protects you from the monsoon. Practical, immediate, real.

\subsection*{The Means}

\begin{versebox}
``All formations are impermanent'' ---\\
``All formations are suffering'' ---\\
``All phenomena are non-self'' ---\\
when one sees this with wisdom,\\
one turns away from suffering.\\
This is the path to purification.
\end{versebox}

That's the means.\footnote{Pali: \textit{``Sabbe saṅkhārā aniccā\ldots sabbe saṅkhārā dukkhā\ldots sabbe dhammā anattā''ti, yadā paññāya passati; / Atha nibbindati dukkhe, esa maggo visuddhiyā''ti.} (Dhp 277--279). The three characteristics of existence (\textit{tilakkhaṇa}) --- impermanence, suffering, and non-self --- are presented here not as philosophical truths to be believed but as \textit{things to be seen with wisdom} (\textit{paññāya passati}). The seeing produces \textit{nibbidā} (disenchantment, turning away) --- not depression or nihilism, but a clear-eyed letting go. The word \textit{visuddhi} (purification) connects to Buddhaghosa's later masterwork, the \textit{Visuddhimagga}.}

The three characteristics. Not as abstract doctrine, but as \textit{method}: see impermanence, see suffering, see non-self. That seeing \textit{is} the path.

\subsection*{The Command}

\begin{versebox}
Like one with good eyes avoiding rough terrain,\\
like a wise person in the living world ---\\
one should avoid what is harmful.
\end{versebox}

That's the command.\footnote{Pali: \textit{``Cakkhumā visamānīva, vijjamāne parakkame; / Paṇḍito jīvalokasmiṃ, pāpāni parivajjaye''ti.} (Ud 4.3). \textit{Cakkhumā} --- ``one with eyes,'' i.e., one who can see. \textit{Visama} --- ``uneven, rough, dangerous.'' The simile is of a sighted person picking their way through treacherous terrain. \textit{Āṇatti} (command, injunction): the Buddha isn't offering a suggestion. He's giving a direct instruction based on clear sight of the terrain.}

Notice the sequence: first acknowledge the pleasure (appeal). Then show the pain (danger). Then show the way out (escape). Then show what happens when you take it (fruit). Then give the specific practice (means). Then --- \textit{do it} (command).

Every teaching the Buddha ever gave follows this pattern, or some subset of it.

\section*{Three Types of Student}

Now Mahākaccāyana reveals something startling: the Buddha doesn't give all six components to everyone. He adapts.\footnote{Pali: \textit{``Tattha bhagavā ugghaṭitaññussa puggalassa nissaraṇaṃ desayati, vipañcitaññussa puggalassa ādīnavañca nissaraṇañca desayati, neyyassa puggalassa assādañca ādīnavañca nissaraṇañca desayati.''} Three types of learner: (1) \textit{ugghaṭitaññū} --- ``one who understands upon a mere hint'' (literally ``one who knows when the topic is merely raised''), (2) \textit{vipañcitaññū} --- ``one who understands when the teaching is elaborated,'' and (3) \textit{neyya} --- ``one who needs to be guided'' (from \textit{neti}, to lead --- the same root as \textit{Netti}). There is a fourth type mentioned in the commentary: \textit{padaparama} --- ``one for whom the words are the utmost'' --- i.e., someone who can memorize the teaching but not realize it in this lifetime. The Netti focuses on the three who can actually benefit.}

\begin{itemize}[itemsep=6pt]
\item The \textbf{quick learner} gets only the escape. They hear ``here's the way out'' and they're already on their feet. They don't need to be convinced. Tell them the solution and they'll take it.

\item The \textbf{elaboration learner} gets the danger \textit{and} the escape. They need to see why the escape matters --- they need to feel the problem before the solution lands.

\item The \textbf{guidable learner} gets the appeal, the danger, \textit{and} the escape. They need the full sequence: yes, this thing is pleasant, yes, it's also dangerous, and here's how to get free. Without acknowledging the appeal first, they won't trust the rest.
\end{itemize}

This is the Buddha as master teacher. He doesn't give a one-size-fits-all lecture. He reads his audience and calibrates.

And \textit{this} is what that stock phrase ``good in the beginning, good in the middle, good in the end'' actually means:\footnote{Pali: \textit{``Soyaṃ dhammavinayo ugghaṭīyanto ugghaṭitaññūpuggalaṃ vineti, tena naṃ āhu `ādikalyāṇo'ti. Vipañcīyanto vipañcitaññūpuggalaṃ vineti, tena naṃ āhu `majjhekalyāṇo'ti. Vitthārīyanto neyyaṃ puggalaṃ vineti, tena naṃ āhu `pariyosānakalyāṇo'ti.''} The Netti reads the three ``levels of goodness'' as three levels of pedagogical elaboration: the opening (\textit{ādi}) is sufficient for quick learners, the middle (\textit{majjhe}) elaboration is sufficient for elaboration learners, and the full conclusion (\textit{pariyosāna}) serves the guidable. This means every sutta is simultaneously three teachings at three levels of detail. A single discourse does triple duty.}

\begin{itemize}[itemsep=4pt]
\item ``Good in the beginning'' --- the opening is enough for the quick learner.
\item ``Good in the middle'' --- the elaboration serves the one who needs more.
\item ``Good in the end'' --- the full treatment guides those who need the whole journey.
\end{itemize}

Every single sutta is simultaneously three teachings for three levels of student. Read that again. One text, three audiences, zero wasted words.

\section*{Four Paths for Four Temperaments}

The Netti then maps this onto four types of practitioner, each with a different temperament, a different dominant faculty, and a different path:\footnote{Pali: \textit{``Tattha catasso paṭipadā, cattāro puggalā. Taṇhācarito mando satindriyena dukkhāya paṭipadāya dandhābhiññāya niyyāti satipaṭṭhānehi nissayehi\ldots''} Four practice-paths (\textit{paṭipadā}) for four persons (\textit{puggalā}), classified by temperament (\textit{carita}) and aptitude (\textit{manda} = slow, \textit{udatta} = keen). The system cross-references temperament (craving vs. views), dominant faculty (mindfulness, concentration, energy, wisdom), path quality (painful vs. pleasant), speed of insight (slow vs. quick), and support base (foundations of mindfulness, jhānas, right efforts, or the four truths). This is one of the most sophisticated personality-based practice models in Buddhist literature.}

\begin{enumerate}[itemsep=6pt]
\item \textbf{Craving temperament, slow learner} --- dominant faculty: mindfulness. Practice is painful, insight comes slowly. Support: the four foundations of mindfulness.

\item \textbf{Craving temperament, keen learner} --- dominant faculty: concentration. Practice is painful, insight comes quickly. Support: the jhānas.

\item \textbf{Views temperament, slow learner} --- dominant faculty: energy. Practice is pleasant, insight comes slowly. Support: the four right efforts.

\item \textbf{Views temperament, keen learner} --- dominant faculty: wisdom. Practice is pleasant, insight comes quickly. Support: the Four Noble Truths.
\end{enumerate}

Both craving-temperament practitioners follow a calm-first path: they develop samatha (tranquility) first, then use it as a base for vipassanā (insight). This leads to liberation of mind --- freedom from passion.\footnote{Pali: \textit{``Ubho taṇhācaritā samathapubbaṅgamāya vipassanā niyyanti rāgavirāgāya cetovimuttiyā.''} --- ``Both craving-temperament persons proceed with insight preceded by calm, toward freedom from passion, toward liberation of mind.'' \textit{Cetovimutti} (liberation of mind) is specifically associated with the overcoming of craving (\textit{rāga}) through calm (\textit{samatha}).}

Both views-temperament practitioners follow an insight-first path: they develop vipassanā first, then deepen it with samatha. This leads to liberation by wisdom --- freedom from ignorance.\footnote{Pali: \textit{``Ubho diṭṭhicaritā vipassanāpubbaṅgame samathena niyyanti avijjāvirāgāya paññāvimuttiyā.''} --- ``Both views-temperament persons proceed with calm preceded by insight, toward freedom from ignorance, toward liberation by wisdom.'' \textit{Paññāvimutti} (liberation by wisdom) is specifically associated with the overcoming of ignorance (\textit{avijjā}) through insight (\textit{vipassanā}).}

And here's the punchline: the calm-first practitioners should be analyzed using the Joyful Round method. The insight-first practitioners should be analyzed using the Lion's Play method.\footnote{Pali: \textit{``Tattha ye samathapubbaṅgamāhi paṭipadāhi niyyanti, te nandiyāvaṭṭena nayena hātabbā, ye vipassanāpubbaṅgamāhi paṭipadāhi niyyanti, te sīhavikkīḷitena nayena hātabbā.''} This is the first time the Netti explicitly connects a specific \textit{hāra} (mode) to a specific \textit{naya} (method). The Teaching Mode doesn't operate in isolation --- it feeds into the larger system. This interlocking of modes and methods is what makes the Netti a system rather than just a list.} The modes and methods aren't separate --- they interlock.

\section*{Three Kinds of Wisdom}

The Netti then maps the teaching process onto three kinds of wisdom:\footnote{Pali: \textit{``Svāyaṃ hāro kattha sambhavati, yassa satthā vā dhammaṃ desayati\ldots so taṃ dhammaṃ sutvā saddhaṃ paṭilabhati.''} The process: hearing the teaching from the Teacher or a respected companion → gaining faith (\textit{saddhā}) → investigation, comparison, and examination → three types of wisdom arise in sequence.}

\begin{enumerate}[itemsep=4pt]
\item \textbf{Wisdom from hearing} --- You hear the teaching, and you examine, compare, and investigate what you've heard. This is the analytical wisdom that comes from study.\footnote{Pali: \textit{``Yā vīmaṃsā ussāhanā tulanā upaparikkhā, ayaṃ sutamayī paññā.''} --- ``The investigation, endeavor, comparison, and examination --- this is wisdom born of hearing.'' \textit{Sutamayī paññā} is the first of the three wisdoms --- it's what you get from listening to, reading, and studying the teachings.}

\item \textbf{Wisdom from reflection} --- You take what you've heard and think it through. You turn it over in your mind, test it, probe it. This is the wisdom that comes from reasoning.\footnote{Pali: \textit{``Tathā sutena nissayena yā vīmaṃsā tulanā upaparikkhā manasānupekkhaṇā, ayaṃ cintāmayī paññā.''} --- ``Using what was heard as a support, the investigation, comparison, examination, and mental reviewing --- this is wisdom born of reflection.'' Note the addition of \textit{manasānupekkhaṇā} (mental reviewing, contemplation) --- reflection adds an element that hearing alone doesn't have.}

\item \textbf{Wisdom from practice} --- When the first two types of wisdom combine with sustained attention, a deeper knowing arises --- either at the level of initial seeing or at the level of full development. This is the wisdom that transforms.\footnote{Pali: \textit{``Imāhi dvīhi paññāhi manasikārasampayuttassa yaṃ ñāṇaṃ uppajjati dassanabhūmiyaṃ vā bhāvanābhūmiyaṃ vā, ayaṃ bhāvanāmayī paññā.''} Two levels: \textit{dassanabhūmi} (``the ground of seeing'' --- initial breakthrough) and \textit{bhāvanābhūmi} (``the ground of development'' --- deepening through sustained practice). \textit{Bhāvanāmayī paññā} isn't something you can get from books or thinking alone --- it requires the combination of study, reflection, \textit{and} sustained meditative attention (\textit{manasikāra}).}
\end{enumerate}

And here's how they map to the three types of student: the quick learner already has both hearing-wisdom and reflection-wisdom. The elaboration learner has hearing-wisdom but not yet reflection-wisdom. The guidable learner has neither --- they need the full treatment from the beginning.\footnote{Pali: \textit{``Yassa imā dve paññā atthi sutamayī cintāmayī ca, ayaṃ ugghaṭitaññū. Yassa sutamayī paññā atthi, cintāmayī natthi, ayaṃ vipañcitaññū. Yassa neva sutamayī paññā atthi na cintāmayī, ayaṃ neyyo.''} This is one of the most psychologically acute observations in the Netti. The quick learner has \textit{both} study and reasoning already developed --- they just need a nudge. The elaboration learner has studied but hasn't thought it through --- they need the teaching developed so they can follow the reasoning. The guidable learner starts fresh --- they need everything from the ground up.}

\section*{The Wheel of Dhamma}

Now the Netti makes its most sweeping claim: the six components of the Teaching Mode map directly onto the Four Noble Truths.\footnote{Pali: \textit{``Sāyaṃ dhammadesanā kiṃ desayati? Cattāri saccāni dukkhaṃ samudayaṃ nirodhaṃ maggaṃ. Ādīnavo ca phalañca dukkhaṃ, assādo samudayo, nissaraṇaṃ nirodho, upāyo āṇatti ca maggo. Imāni cattāri saccāni. Idaṃ dhammacakkaṃ.''} The mapping: danger + fruit = suffering (the First Truth), appeal = origin (the Second Truth --- craving for pleasure is the origin of suffering), escape = cessation (the Third Truth), and means + command = the path (the Fourth Truth). Then the stunning conclusion: \textit{``Idaṃ dhammacakkaṃ''} --- ``This is the Wheel of Dhamma.'' Every teaching that follows the six-component structure is a turning of the Wheel.}

\begin{itemize}[itemsep=4pt]
\item Danger and fruit = \textbf{Suffering} --- the First Noble Truth
\item Appeal = \textbf{Origin} --- the Second Noble Truth (craving for pleasure is the origin of suffering)
\item Escape = \textbf{Cessation} --- the Third Noble Truth
\item Means and command = \textbf{The Path} --- the Fourth Noble Truth
\end{itemize}

``This is the Wheel of Dhamma.''

Every teaching that follows this six-part pattern --- appeal, danger, escape, fruit, means, command --- is a turning of the Wheel. Not just the famous first sermon at Deer Park. \textit{Every} teaching. The Wheel turns every time the pattern plays out.

And then the Netti quotes the Buddha himself: ``This is suffering --- at Bārāṇasī, in the Deer Park at Isipatana, the unsurpassed Wheel of Dhamma was set rolling, irreversible by any ascetic or brahmin or god or Māra or Brahmā or anyone in the world.''\footnote{This is the stock passage from SN 56.11, the \textit{Dhammacakkappavattana Sutta} --- the first discourse. The Netti's point is that the Teaching Mode doesn't just apply to this one sutta. The structure of the first sermon is the structure of \textit{all} the Buddha's teaching. The Wheel didn't turn once at Deer Park --- it turns in every sutta.}

The implication is staggering. The Canon isn't a collection of separate teachings. It's a single Wheel, turning the same pattern over and over, in text after text, adapted to context after context. And the Teaching Mode is how you see it.

\ornament

\noindent That's the first mode, fully demonstrated. In the next chapter, we'll see the Investigation Mode --- how to break any sutta down to its question, its answer, and its summary --- applied to one of the most beautiful question-and-answer exchanges in the entire Canon.
